\documentclass[11pt]{article}
\usepackage[a4paper,margin=1in]{geometry}
\usepackage{subfiles}
\usepackage{amsmath,amssymb,bm}
\usepackage{hyperref}
\usepackage{booktabs}
\usepackage{array}
\usepackage{mathtools}

\title{Appendix 2 (to main theory): Scalar Energy--Momentum Tensor and Conservation Identities in One--Metric ISPG}
\author{}
\date{}

\begin{document}
\let\phi\varphi
\maketitle

\section*{Purpose and scope}
This appendix records the explicit energy--momentum tensor associated with the main theory scalar sector and the corresponding conservation identities. The intent is to provide a single, citable location for:
(i) the explicit form of $T^{(\phi)}_{\mu\nu}$ used in the Einstein equations,
(ii) the scalar equation of motion as a conservation identity,
and (iii) the covariant conservation of matter stress--energy under minimal coupling.

\section*{Non-negotiable consistency with the main theory (one-metric formulation)}
We work strictly within the main theory posture:
\begin{itemize}
\item matter is minimally coupled to the single physical metric $g_{\mu\nu}$ (no disformal mapping and no frame change);
\item the covariant backbone is
\begin{equation}
S=\frac{1}{16\pi G}\int d^4x\sqrt{-g}\left[R+\frac12\,g^{\mu\nu}\partial_\mu\phi\,\partial_\nu\phi\right]+S_m[g_{\mu\nu},\psi],
\label{eq:action_app2}
\end{equation}
with the same metric $g_{\mu\nu}$ as used elsewhere in the main theory package.
\end{itemize}

\section{Metric variation and field equations}\label{sec:metricvar}
Varying \eqref{eq:action_app2} with respect to $g^{\mu\nu}$ gives
\begin{equation}
G_{\mu\nu} + S_{\mu\nu} = 8\pi G\,T^{(m)}_{\mu\nu},
\label{eq:einsteinS_app2}
\end{equation}

where $T^{(m)}_{\mu\nu}$ is defined in the standard way by
\begin{equation}
\delta S_m = -\frac12\int d^4x\sqrt{-g}\,T^{(m)}_{\mu\nu}\,\delta g^{\mu\nu}.
\label{eq:Tmdef_app2}
\end{equation}
The explicit algebraic form of $S_{\mu\nu}$ and its relation to the canonical scalar tensor $T^{(\phi)}_{\mu\nu}$ are given in Section~1.1 below.
The tensor $S_{\mu\nu}$ is the contribution coming from the scalar kinetic term inside the $1/(16\pi G)$ prefactor.

\subsection{Two equivalent (bookkeeping) conventions}
It is often convenient to rewrite \eqref{eq:einsteinS_app2} in the more familiar Einstein form
\begin{equation}
G_{\mu\nu}=8\pi G\,T^{(m)}_{\mu\nu}+\kappa\,T^{(\phi)}_{\mu\nu},
\label{eq:einsteinT_app2}
\end{equation}
with a purely conventional constant $\kappa$ that tracks
the sign choice.
\textbf{Notation.}  This appendix denotes the scalar
energy-momentum tensor by $T^{(\phi)}_{\mu\nu}$; the
main text uses $\tau^{(\varphi)}_{\mu\nu}$ for the same
object (with the identification $\phi \equiv \varphi$).
A clean choice is to define $T^{(\phi)}_{\mu\nu}$ in
the \emph{canonical} algebraic form
\begin{equation}
T^{(\phi)}_{\mu\nu}\equiv \partial_\mu\phi\,\partial_\nu\phi-\frac12 g_{\mu\nu}(\partial\phi)^2,
\qquad (\partial\phi)^2\equiv g^{\alpha\beta}\partial_\alpha\phi\,\partial_\beta\phi,
\label{eq:Tphi_canonical}
\end{equation}
in which case the main theory sign choice corresponds to
\begin{equation}
\kappa=-\frac12.
\label{eq:kappa}
\end{equation}
Equivalently, one may absorb $\kappa$ into the definition and write directly
\begin{equation}
S_{\mu\nu}= -\kappa\,T^{(\phi)}_{\mu\nu}
\quad\Longleftrightarrow\quad
S_{\mu\nu}=\frac12\left(\partial_\mu\phi\,\partial_\nu\phi-\frac12 g_{\mu\nu}(\partial\phi)^2\right),
\label{eq:Smunu_app2}
\end{equation}
so that \eqref{eq:einsteinS_app2} and \eqref{eq:einsteinT_app2} are identical statements.
The physics is unchanged by this bookkeeping choice; it simply makes the sign tracking explicit in a form consistent across the main theory appendices.

\section{Scalar variation and equation of motion}\label{sec:scalarvar}
Varying \eqref{eq:action_app2} with respect to $\phi$ under minimal coupling ($\delta S_m/\delta\phi\equiv 0$) yields
\begin{equation}
\square\phi = 0 \qquad \text{(vacuum and matter regions alike)} ,
\label{eq:boxphi_app2}
\end{equation}
The reduced-sector consistency relation used throughout the main theory,
\begin{equation}
\square\phi = -\frac{8\pi G}{c^4}\,T,
\qquad T\equiv g^{\mu\nu}T^{(m)}_{\mu\nu},
\end{equation}
is \emph{not} obtained from this direct $\phi$-variation; it is the trace consistency of the metric equation \eqref{eq:einsteinT_app2} after substituting the bi--conformal ansatz $g_{\mu\nu}[\phi]$ (Appendix~3 and Appendix~1, §Scalar field equation). For the purposes of this appendix, the key identity is the covariant relation between the divergence of $T^{(\phi)}_{\mu\nu}$ and the scalar equation:
\begin{equation}
\nabla^\mu T^{(\phi)}_{\mu\nu} = (\square\phi)\,\partial_\nu\phi.
\label{eq:divTphi_app2}
\end{equation}
Thus, on-shell ($\square\phi=0$ in vacuum), the scalar energy--momentum is covariantly conserved.

\section{Conservation identities and Bianchi structure}\label{sec:cons_app2}
Taking the covariant divergence of \eqref{eq:einsteinT_app2} and using the Bianchi identity $\nabla^\mu G_{\mu\nu}=0$ gives
\begin{equation}
\nabla^\mu T^{(m)}_{\mu\nu} = -\frac{\kappa}{8\pi G}\,\nabla^\mu T^{(\phi)}_{\mu\nu}.
\label{eq:balance_app2}
\end{equation}
Because matter is minimally coupled to $g_{\mu\nu}$, its stress--energy is conserved by construction,
\begin{equation}
\nabla^\mu T^{(m)}_{\mu\nu}=0,
\label{eq:divTm}
\end{equation}
and therefore \eqref{eq:balance_app2} reduces to $\nabla^\mu T^{(\phi)}_{\mu\nu}=0$ whenever the scalar equation is satisfied (via \eqref{eq:divTphi_app2}). This closes the consistency loop: metric equations + scalar equation + minimal coupling imply consistent covariant conservation.

\section{Useful algebraic identities}\label{sec:ids}
For later reference, the trace of the canonical-form tensor \eqref{eq:Tphi_canonical} is
\begin{equation}
T^{(\phi)}\equiv g^{\mu\nu}T^{(\phi)}_{\mu\nu}= -(\partial\phi)^2,
\end{equation}
and for a static configuration where $\phi=\phi(r)$ one has
\begin{equation}
(\partial\phi)^2=g^{rr}(\partial_r\phi)^2,
\qquad
T^{(\phi)}_{tt}=-\frac12 g_{tt}(\partial\phi)^2,
\qquad
T^{(\phi)}_{rr}=(\partial_r\phi)^2-\frac12 g_{rr}(\partial\phi)^2,
\end{equation}
with angular components following by symmetry.
For a purely radial static configuration, substituting
$(\partial\phi)^2=g^{rr}(\partial_r\phi)^2$ gives
$T^{(\phi)}_{rr}=\frac12(\partial_r\phi)^2$,
which is \emph{not} zero.  The mixed components are
\begin{equation}
T^{(\phi)\,t}{}_{t}
  =g^{tt}T^{(\phi)}_{tt}=-\tfrac12 g^{rr}(\partial_r\phi)^2,
\qquad
T^{(\phi)\,r}{}_{r}
  =g^{rr}T^{(\phi)}_{rr}=+\tfrac12 g^{rr}(\partial_r\phi)^2,
\end{equation}
so $T^{(\phi)\,t}{}_{t}=-T^{(\phi)\,r}{}_{r}$
and therefore
\begin{equation}\label{eq:sum_identity}
T^{(\phi)\,t}{}_{t}+T^{(\phi)\,r}{}_{r}=0.
\end{equation}
Via the vacuum Einstein equation this forces
$G^{t}{}_{t}+G^{r}{}_{r}=0$,
which---for a general static isotropic metric
$ds^{2}=-e^{2\alpha}c^{2}dt^{2}+e^{2\beta}(dr^{2}+r^{2}d\Omega^{2})$---implies
$(\alpha+\beta)'=0$.
Asymptotic flatness ($\alpha,\beta\to0$) then gives
$\alpha=-\beta$, hence $g_{tt}\cdot g_{rr}=-c^{2}$.
This is the derivation referenced in the main theory.
\section{Derived vs.\ stated: compact ledger}\label{sec:ledger_app2}
\begin{center}
\begin{tabular}{@{}p{0.28\linewidth}p{0.62\linewidth}@{}}
\toprule
\textbf{Derived here} &
Explicit scalar energy--momentum tensor \eqref{eq:Tphi_canonical} and its divergence identity \eqref{eq:divTphi_app2}; bookkeeping relation to the main theory Einstein equation via $\kappa$ \eqref{eq:kappa}; conservation loop \eqref{eq:balance_app2}--\eqref{eq:divTm}.\\
\midrule
\textbf{Used elsewhere in the package} &
Weak-field PPN/deflection/Shapiro appendices rely on $\Phi=\Psi$ and consistent sourcing; Appendix~11 uses the constraint/hyperbolicity analysis together with the conservation structure; GW appendix uses minimal coupling and the one-metric posture.\\
\bottomrule
\end{tabular}
\end{center}

\section*{Takeaway}
In the main theory one-metric formulation, the scalar sector contributes an explicit stress--energy tensor. In canonical algebraic form it is
\[
T^{(\phi)}_{\mu\nu}=\partial_\mu\phi\,\partial_\nu\phi-\frac12 g_{\mu\nu}(\partial\phi)^2,
\]
and the main theory sign bookkeeping corresponds to $\kappa=-1/2$ in $G_{\mu\nu}=8\pi G T^{(m)}_{\mu\nu}+\kappa T^{(\phi)}_{\mu\nu}$. Minimal coupling implies $\nabla^\mu T^{(m)}_{\mu\nu}=0$, while $\nabla^\mu T^{(\phi)}_{\mu\nu}=(\square\phi)\partial_\nu\phi$ closes the conservation identities on-shell.

\end{document}
